\section{Materiale di consegna e riproducibilità}
\label{sec:deliverables}

\subsection{Esecuzione}

Il punto di ingresso è \texttt{src/main.py}. La variabile \texttt{MODE}
seleziona tre operazioni distinte:
\begin{itemize}
    \item \texttt{generate}: crea il dataset con i seed e il budget
    dell'insegnante specificati nella configurazione;
    \item \texttt{train}: legge il dataset, prepara le feature, separa
    le istanze, addestra e salva i parametri del modello;
    \item \texttt{compare}: legge dataset e modello esistenti, calcola
    i riferimenti e confronta le ricerche sui budget scelti.
\end{itemize}
Dalla radice del sottoprogetto si esegue \texttt{python3 -m src.main}.
La configurazione corrente usa \texttt{START\_PATH = (0,)}, le istanze
$50$ e $51$ e i quattro budget della tabella. I risultati, inclusi
\texttt{optimal\_cost}, \texttt{gap\_percent} e le statistiche di ciascuna
ricerca, sono salvati in \path{results/search_comparison.json}.
Le esecuzioni riportate utilizzano Python $3.14.6$ e
\texttt{scikit-learn} $1.9.1$. La dipendenza è elencata in
\texttt{requirements.txt}.

\subsection{Codice, dati e risultati}

Il programma contiene la costruzione delle istanze e dei target, il codice
per addestrare il predittore e il confronto sperimentale A/B. I file
necessari a riprodurre i risultati sono organizzati come segue:
\begin{itemize}
    \item \path{src/algorithms/baseline.py}: ricerca DFS con B\&B;
    \item \path{src/algorithms/predictor.py}: feature, addestramento,
    lettura dei parametri e previsione;
    \item \path{src/algorithms/predictive_search.py}: ordinamento predittivo
    dei successori e ricerca a budget limitato;
    \item \path{src/experiments/generate_dataset.py}: istanze riproducibili
    e produzione dei target tramite ricerca insegnante;
    \item \path{data/examples.json}, \path{data/result.json} e
    \path{data/model.json}: dati, feature e parametri appresi;
    \item \path{results/search_comparison.json}: costi, percorsi, riferimenti
    e statistiche delle sedici esecuzioni.
\end{itemize}
Il materiale consente di riprodurre le sedici esecuzioni A/B riportate
nella relazione. Il registro delle interazioni AI è riportato
nell'Allegato~A.
